\documentclass[10pt]{report}
\usepackage[utf8]{inputenc}

\input{general_preamble}

% DEFINE TITLE STYLES
\makeatletter 
\renewcommand\maketitle{
{\begin{center}
{\Large \bfseries \@title }\\
{\small \emph{First created}: \DTMdisplaydate{2026}{01}{22}{-1}\hfill%
\setstretch{1.0} \small \hfill \emph{Last update}: \today}
%{\Large  \@author}\\[4ex] 
\end{center}}} % Note the extra }
\makeatother

%-------------------------------------------------------------------------

\begin{document}

%\sisetup{parse-numbers = false}
\DTMsetdatestyle{cvdateformat}

\title{\vspace{-3em} Multiplexed mRNA \textit{in-situ} Hybridization Chain Reaction (HCR) in Placozoa}
\maketitle
\thispagestyle{plain}
\markboth{HCR in Placozoa}{}

\section*{\vspace{-1em} \textendash{} Protocol \textendash{}}
\markright{Protocol}

\subsection*{Sample fixation}

\begin{enumerate}[series = steps]
	\item Gently transfer Placozoa with a wide-boar pipette tip onto a small petri dish or a multi-well plate, and allow the animals to settle for \qtyrange{15}{30}{\minute}.
	\item Remove as much culture water as possible, taking care of not exposing animals to air.
	\item \step{fix}{\qty{> 12}{\hour} (overnight)}{the required amount}{fixative solution}{\fourdegree}
\end{enumerate}

\bigskip\alert{Prepare fresh fixative solution soon before usage.}

\begin{enumerate}[resume = steps]
	\item Gently transfer fixed samples into a \qty{1.5}{\ml} low-bind tube.
	\item \step{wash}{\qtytimes{3}{10}{\minute}}{\fivehunmicrol{} (or \onemil)}{\pbs}{RT}
	\item \step{dehydrate}{\qtytimes{4}{10}{\minute}}{\fivehunmicrol{} (or \onemil)}{ice-cold methanol}{RT}
	\item Store samples at \minustwenty{} until usage.
\end{enumerate}

\subsection*{Day 1 \textendash{} Sample preparation}

\alert{In each step, samples can be rocked on a nutator or on an orbital (horizontal) shaker.}
\alert{Adding and removing solutions should be performed under a dissecting scope to avoid losing a large number of embryos.}

\begin{enumerate}[series = steps]
	\item Transfer \qtyrange{\sim 5}{10}{Placozoa} into each well of a 96 well plate.
	\item Rehydrate samples for \underline{\qty{5}{\minute}} each with a graded series of \qty{250}{\ul} of \textbf{methanol (\qtylist{75;50;25}{\percent}) in PBST} at RT.
	\item \step{wash}{\qty{5}{\minute}}{\qty{250}{\ul}}{PBST}{RT}
	\item \step{rock}{\tenmin}{\qty{250}{\ul}}{\qty{10}{\ug\per\ml} proteinase K}{\qty{37}{\degreeCelsius}}
	\item \step{wash}{two times without incubation}{\qty{250}{\ul}}{PBST}{RT}
	\item \step{post-fix}{\quarter}{\qty{250}{\ul}}{\qty{4}{\percent} PFA}{RT}
\end{enumerate}

\bigskip\alert{Pre-heat probe-hybridization buffer to \thirtysevendegree{} (you will need \qty{200}{\ul} per tube [\qty{100}{\ul} for the pre-hybridization, and \qty{100}{\ul} for the probe solution]).}

\begin{enumerate}[resume = steps]
	\item \step{wash}{\qtytimes{4}{5}{\minute}}{\qty{250}{\ul}}{PBST}{RT}
\end{enumerate}

\subsection*{Day 1 \textendash{} HCR detection stage}
\begin{enumerate}[series = steps]
	\item \step{pre-hybridize}{\halfhour}{\qty{100}{\ul}}{pre-warmed probe hybridization buffer}{\thirtysevendegree}
\end{enumerate}

\bigskip\alert{If reusing probes, pre-warm probe solution to \thirtysevendegree, then skip to Step 3.}
\alert{At this step, samples can be stored in probe hybridization buffer at \minustwenty{} for several months.}

\begin{enumerate}[resume = steps]
	\item Prepare probe solution by adding \qty{1}{\ul} (\qty{1}{\pmol}) of each \textbf{probe set \qty{1}{\micro\molar} stock solution} to \qty{100}{\ul} of \textbf{probe hybridization buffer} at \thirtysevendegree.
	\item Add probe solution to the well to reach a final volume of \qty{200}{\ul} (final concentration of probes of \qty{5}{\nano\molar}).
	\item Seal the plate with film and incubate samples for \underline{\qty{> 12}{\hour} (overnight)} at \thirtysevendegree.
\end{enumerate}

\bigskip\alert{Pre-heat probe wash buffer at \thirtysevendegree{} (you will need \qty{250}{\ul} per tube per 4 washes).}

\subsection*{Day 2 \textendash{} HCR amplification stage}

\alert{Pre-equilibrate amplification buffer to room temperature (you will need \qty{100}{\ul} per tube [\qty{50}{\ul} for pre-amplification, and \qty{50}{\ul} for the hairpin solution]).}

\begin{enumerate}[series = steps]
	\item Add \qty{\sim 150}{\ul} of probe wash buffer to the sample and incubate for \tenmin.
	\item \step{wash}{\qtytimes{4}{15}{\minute}}{\qty{250}{\ul}}{probe wash buffer}{\thirtysevendegree}
\end{enumerate}

\bigskip\alert{Probe solution can be saved and reused for 2--3 times; store at \minustwenty.}
\alert{If probe solution is too dense and does not allow a safe removal without loosing samples, first add \qty{\sim 150}{\ul} of pre-warmed probe wash buffer, then proceed with step 1.}

\begin{enumerate}[resume = steps]
	\item \step{wash}{\qtytimes{2}{5}{\minute}}{\qty{250}{\ul}}{\ssct}{RT}
	\item \step{pre-amplify}{\halfhour}{\qty{50}{\ul}}{amplification buffer}{RT}
	\item Snap cool in separate tubes \qty{2}{\ul} of \textbf{hairpin H1} (\qty{6}{\pmol}) and \qty{2}{\ul} of \textbf{hairpin H2} (\qty{6}{\pmol}) \textbf{\qty{3}{\micro\molar} stock solutions}: heat at \qty{95}{\degreeCelsius} for \qty{90}{\s} and cool to RT in the dark for \halfhour.
	      %\note{doubling hairpin concentration can help to boost signal.}
\end{enumerate}

\bigskip\alert{If reusing hairpin solutions, heat at \qty{95}{\degreeCelsius} and cool at RT, then skip at Step 6.}

\begin{enumerate}[resume = steps]
	\item Prepare hairpin solution by adding \textbf{snap-cooled H1 hairpins} and \textbf{snap-cooled H2 hairpins} to \qty{100}{\ul} of \textbf{amplification buffer} at RT.
	\item Add hairpin solution to the well to reach a final volume of \qty{100}{\ul} (final concentration of each hairpin of \qty{60}{\nano\molar}).
	\item Seal the plate with film and incubate samples for \underline{\qty{> 12}{\hour} (overnight)} at RT.
\end{enumerate}

\subsection*{Day 3 \textendash{} HCR conclusion}

\begin{enumerate}[series = steps]
	\item Add \qty{\sim 150}{\ul} of \ssct to the sample and incubate for \tenmin.
	\item \step{wash}{\qtytimes{2}{5}{\minute} and \qtytimes{2}{30}{\minute}}{\qty{250}{\ul}}{\ssct}{RT}
\end{enumerate}

\bigskip\alert{dsDNA can be stained on the second \qty{30}{\minute} wash by using a 1:500 DAPI solution.}
\alert{Hairpin solution can be saved and reused for 2--3 times; store at \minustwenty.}
\alert{At this step, samples can be stored in the dark in \ssct for several days at room temperature.}

\subsection*{Day 3 \textendash{} Sample mounting}

\begin{enumerate}[series = steps]
	\item \step{incubate}{\qtyrange{30}{60}{\minute}}{\qty{250}{\ul}}{\qty{50}{\percent} glycerol in \pbs}{RT}
	\item \step{incubate}{\qtyrange{30}{60}{\minute}}{\qty{\sim 250}{\ul}}{anti-fade mounting medium}{RT}
	\item Collect samples and place them on a cleaned (bridged, if necessary) slide.
\end{enumerate}

\bigskip\alert{Adjust the amount of mounting medium according to the cover glass dimensions, the quantity of samples, and the thickness of the bridge.}

\begin{enumerate}[resume = steps]
	\item	Seal the slide with nail polish and store in the dark at \fourdegree.
\end{enumerate}

\clearpage

\section*{\textendash{} Recipes \textendash{}}
\setstretch{0.95}
\markright{Recipes}

\subsection*{10\per{} PHEM buffer}

\begin{table}[H]
	\centering
	\begin{tabular}{r
		S[table-format = 4.1, exponent-mode = fixed, fixed-exponent = 0, round-mode = places, round-precision = 1]@{\,} % disable scientific notation
		l
		S[table-format = 3.2, table-align-text-post = false, exponent-mode = fixed, fixed-exponent = 0, round-mode = places, round-precision = 2] % disable scientific notation
		c
		}
		\textbf{Compound}				& \multicolumn{2}{c}{\textbf{Quantity}}	& \multicolumn{1}{c}{\textbf{\begin{tabular}[c]{@{}c@{}}Molar mass\\ (\unit{\g\per\mole})\end{tabular}}}	& \multicolumn{1}{l}{\textbf{\begin{tabular}[c]{@{}c@{}}Final concentration\\ (for \qty{50}{\ml})\end{tabular}}}	\\*[0.4cm]
		PIPES							& 9.07	& \unit{\g}						& 302.37																									& \qty{600}{\milli\molar}																							\\
		HEPES							& 2.98	& \unit{\g}						& 238.3                                                                                     				& \qty{250}{\milli\molar}              																				\\
		\qty{0.5}{\molar} EGTA pH 0.8	& 10	& \unit{\ml}					& 380.35                                                                                     				& \qty{100}{\milli\molar}              																				\\
		\qty{1}{\molar} \ce{MgCl2}      & 1		& \unit{\ml}					& 95.211                                                                                     				& \qty{20}{\milli\molar}              																				\\
		Ultrapure water					& \multicolumn{2}{c}{\NA}               & \NA                                                                                                    	& \NA																												\\
	\end{tabular}
\end{table}

\begin{enumerate}
	\item Combine \qty{9.07}{\g} of PIPES and \qty{2.98}{\g} of HEPES into a becker.
	\item Add \qty{25}{\ml} of ultrapure water to the becker and start stirring gently with a magnetic stirrer.
	\item Add \qty{10}{\ml} of \qty{0.5}{\molar} EGTA pH 0.8 and \qty{1}{\ml} of \qty{1}{\molar} \ce{MgCl2}.
	\item Gently raise the pH of the solution to 7.5 by adding \ce{KOH} to allow full solubilization of components.
	\item Transfer the solution to a graduated cylinder and fill up to \qty{50}{\ml} with ultrapure water.
	\item Filter the solution through a \qty{0.22}{\um} filter. Autoclave if required.
\end{enumerate}


\subsection*{Fixative solution}

\begin{table}[H]
	\centering
	\begin{tabular}{r
		S[table-format = 4.1, exponent-mode = fixed, fixed-exponent = 0, round-mode = places, round-precision = 1]@{\,} % disable scientific notation
		l
		S[table-format = 3.2, table-align-text-post = false, exponent-mode = fixed, fixed-exponent = 0, round-mode = places, round-precision = 2] % disable scientific notation
		c
		}
		\textbf{Compound}						& \multicolumn{2}{c}{\textbf{Quantity}}	& \multicolumn{1}{c}{\textbf{\begin{tabular}[c]{@{}c@{}}Molar mass\\ (\unit{\g\per\mole})\end{tabular}}}	& \multicolumn{1}{l}{\textbf{\begin{tabular}[c]{@{}c@{}}Final concentration\\ (for \qty{20}{\ml})\end{tabular}}}	\\*[0.4cm]
		10\per{} PHEM buffer					& 3		& \unit{\ml}					& \NA																										& 1.5\per{}																											\\
		\qty{37}{\percent}	(para)formaldehyde	& 2.16	& \unit{\ml}					& 30.031                                                                                     				& \qty{4}{\percent}              																					\\
		Ultrapure water							& \multicolumn{2}{c}{\NA}               & \NA                                                                                                    	& \NA																												\\
	\end{tabular}
\end{table}

\begin{enumerate}
	\item Combine \qty{3}{\ml} of 10\per{} PHEM buffer and \qty{2.16}{\ml} of \qty{37}{\percent} (para)formaldehyde into a graduated cylinder.
	\item Fill up to \qty{20}{\ml} with ultrapure water.
	\item Use fresh.	
\end{enumerate}



\subsection*{\pbs{} with \qty{0.1}{\percent} Tween 20 (PBST or PTw)}

\begin{table}[H]
	\centering
	\begin{tabular}{r
		S[table-format = 4.1, exponent-mode = fixed, fixed-exponent = 0, round-mode = places, round-precision = 1]@{\,} % disable scientific notation
		l
		S[table-format = 3.2, table-align-text-post = false, exponent-mode = fixed, fixed-exponent = 0, round-mode = places, round-precision = 2] % disable scientific notation
		c
		}
		\textbf{Compound} & \multicolumn{2}{c}{\textbf{Quantity}} & \multicolumn{1}{c}{\textbf{\begin{tabular}[c]{@{}c@{}}Molar mass\\ (\unit{\g\per\mole})\end{tabular}}} & \multicolumn{1}{l}{\textbf{\begin{tabular}[c]{@{}c@{}}Final concentration\\ (for \qty{50}{\ml})\end{tabular}}}                        \\*[0.4cm]
		Tween 20          & 50                                    & \unit{\ul}                                                                                             & 1227.54                                                                                                        & \qty{0.1}{\percent} \\
		20\per{} PBS      & 5                                     & \unit{\ml}                                                                                             & \NA                                                                                                            & 1\per{}              \\
		Ultrapure water   & \multicolumn{2}{c}{\NA}               & \NA                                                                                                    & \NA
	\end{tabular}
\end{table}

\begin{enumerate}
	\item Add \qty{5}{\ml} of 20\per{} PBS to a graduated cylinder.
	\item Add \qty{50}{\ul} of Tween 20.
	\item Fill up to \qty{50}{\ml} with ultrapure water.
	\item Filter the solution through a \qty{0.22}{\um} filter. Autoclave if required.
\end{enumerate}

\subsection*{5\per{} SSC \qty{0.1}{\percent} Tween 20 (SSCT)}

\begin{table}[H]
	\centering
	\begin{tabular}{r
		S[table-format = 4.1, exponent-mode = fixed, fixed-exponent = 0, round-mode = places, round-precision = 1]@{\,} % disable scientific notation
		l
		S[table-format = 3.2, table-align-text-post = false, exponent-mode = fixed, fixed-exponent = 0, round-mode = places, round-precision = 2] % disable scientific notation
		c
		}
		\textbf{Compound} & \multicolumn{2}{c}{\textbf{Quantity}} & \multicolumn{1}{c}{\textbf{\begin{tabular}[c]{@{}c@{}}Molar mass\\ (\unit{\g\per\mole})\end{tabular}}} & \multicolumn{1}{l}{\textbf{\begin{tabular}[c]{@{}c@{}}Final concentration\\ (for \qty{50}{\ml})\end{tabular}}}                        \\*[0.4cm]
		Tween 20          & 50                                    & \unit{\ul}                                                                                             & 1227.54                                                                                                        & \qty{0.1}{\percent} \\
		20\per{} SSC      & 12.5                                  & \unit{\ml}                                                                                             & \NA                                                                                                            & 5\per{}              \\
		Ultrapure water   & \multicolumn{2}{c}{\NA}               & \NA                                                                                                    & \NA
	\end{tabular}
\end{table}

\begin{enumerate}
	\item Add \qty{12.5}{\ml} of 20\per{} SSC to a graduated cylinder.
	\item Add \qty{50}{\ul} of Tween 20.
	\item Fill up to \qty{50}{\ml} with ultrapure water.
	\item Filter the solution through a \qty{0.22}{\um} filter. Autoclave if required.
\end{enumerate}

\subsection*{Proteinase K working solution}

\begin{table}[H]
	\centering
	\begin{tabular}{r
		S[table-format = 4.2, exponent-mode = fixed, fixed-exponent = 0, round-mode = places, round-precision = 2]@{\,} % disable scientific notation
		l
		S[table-format = 3.2, table-align-text-post = false, exponent-mode = fixed, fixed-exponent = 0, round-mode = places, round-precision = 2] % disable scientific notation
		c
		}
		\textbf{Compound}                                                                                              & \multicolumn{2}{c}{\textbf{Quantity}} & \multicolumn{1}{c}{\textbf{\begin{tabular}[c]{@{}c@{}}Molar mass\\ (\unit{\g\per\mole})\end{tabular}}} & \multicolumn{1}{l}{\textbf{\begin{tabular}[c]{@{}c@{}}Final concentration\\ (for \qty{10}{\ml})\end{tabular}}}                         \\*[0.4cm]
		\multicolumn{1}{c}{\begin{tabular}[c]{@{}r@{}}\qty{20}{\mg\per\ml} proteinase K\\ stock solution\end{tabular}} & 5.0                                   & \unit{\ul}                                                                                             & \NA                                                                                                            & \qty{10}{\ug\per\ml} \\*[0.4cm]
		PBST                                                                                                           & \multicolumn{2}{c}{\NA}               & \NA                                                                                                    & \NA
	\end{tabular}
\end{table}

\begin{enumerate}
	\item Add \qty{5}{\ul} of \qty{20}{\mg\per\ml} proteinase K stock solution to \qty{8}{\ml} of PBST.
	\item Fill up to \qty{10}{\ml} with PBST.
\end{enumerate}

\subsection*{Urea-based probe hybridization buffer}

\begin{table}[H]
	\centering
	\begin{tabular}{r
		S[table-format = 4.2, exponent-mode = fixed, fixed-exponent = 0, round-mode = places, round-precision = 2]@{\,} % disable scientific notation
		l
		S[table-format = 3.2, table-align-text-post = false, exponent-mode = fixed, fixed-exponent = 0, round-mode = places, round-precision = 2] % disable scientific notation
		c
		}
		\textbf{Compound}                                                                                              & \multicolumn{2}{c}{\textbf{Quantity}} 						& \multicolumn{1}{c}{\textbf{\begin{tabular}[c]{@{}c@{}}Molar mass\\ (\unit{\g\per\mole})\end{tabular}}} & \multicolumn{1}{l}{\textbf{\begin{tabular}[c]{@{}c@{}}Final concentration\\ (for \qty{10}{\ml})\end{tabular}}}                         \\*[0.4cm]
		Urea																										   & 2.4                                   & \unit{\g}			& 60.06                                                                                            		 & \qty{4}{\molar} 																														  \\
		20\per{} SSC                                                                                                   & 2.5                                   & \unit{\ml}         & \NA{}                                                                                  				 & 5\per{}																																  \\
		\qty{1}{\molar} citric acid																			           & 90									   & \unit{\ul}			& 192.124																								 & \qty{9}{\milli\molar}																												  \\
		Tween-20																									   & 10									   & \unit{\ul}			& 1227.54																							     & \qty{0.1}{\percent}                                                                                                                    \\
		50\per{} Denhardts's solution																				   & 200								   & \unit{\ul}			& \NA{}																									 & 1\per{}																																  \\
		\qty{50}{\percent} dextran sulfate																	           & 2   								   & \unit{\ml}			& \NA{}																									 & \qty{10}{\percent}																													  \\
		\qty{10}{\mg\per\ml} heparin																				   & 50									   & \unit{\ul}			& \NA{}																									 & \qty{50}{\ug\per\ml}																													  \\
	\end{tabular}
\end{table}

\begin{enumerate}
	\item Add \qty{2.4}{\g} of urea to a becker.
	\item Combine all the other components fresh.
	\item Stir thoroughly to allow complete solubilization of urea.
	\item Fill up to \qty{10}{\ml} with ultrapure water.
\end{enumerate}

\subsection*{Urea-based probe wash buffer}

\begin{table}[H]
	\centering
	\begin{tabular}{r
		S[table-format = 4.2, exponent-mode = fixed, fixed-exponent = 0, round-mode = places, round-precision = 2]@{\,} % disable scientific notation
		l
		S[table-format = 3.2, table-align-text-post = false, exponent-mode = fixed, fixed-exponent = 0, round-mode = places, round-precision = 2] % disable scientific notation
		c
		}
		\textbf{Compound}                                                                                              & \multicolumn{2}{c}{\textbf{Quantity}} 						& \multicolumn{1}{c}{\textbf{\begin{tabular}[c]{@{}c@{}}Molar mass\\ (\unit{\g\per\mole})\end{tabular}}} & \multicolumn{1}{l}{\textbf{\begin{tabular}[c]{@{}c@{}}Final concentration\\ (for \qty{10}{\ml})\end{tabular}}}                         \\*[0.4cm]
		Urea																										   & 2.4                                   & \unit{\g}			& 60.06                                                                                            		 & \qty{4}{\molar} 																														  \\
		20\per{} SSC                                                                                                   & 2.5                                   & \unit{\ml}         & \NA{}                                                                                  				 & 5\per{}																																  \\
		\qty{1}{\molar} citric acid																			           & 90									   & \unit{\ul}			& 192.124																								 & \qty{9}{\milli\molar}																												  \\
		Tween-20																									   & 10									   & \unit{\ul}			& 1227.54																							     & \qty{0.1}{\percent}                                                                                                                    \\
		\qty{10}{\mg\per\ml} heparin																				   & 50									   & \unit{\ul}			& \NA{}																									 & \qty{50}{\ug\per\ml}																													  \\
	\end{tabular}
\end{table}

\begin{enumerate}
	\item Add \qty{2.4}{\g} of urea to a becker.
	\item Combine all the other components fresh.
	\item Stir thoroughly to allow complete solubilization of urea.
	\item Fill up to \qty{10}{\ml} with ultrapure water.
\end{enumerate}


\subsection*{1:500 DAPI staining solution}

\begin{table}[H]
	\centering
	\begin{tabular}{r
		S[table-format = 4.1, exponent-mode = fixed, fixed-exponent = 0, round-mode = places, round-precision = 1]@{\,} % disable scientific notation
		l
		S[table-format = 3.2, table-align-text-post = false, exponent-mode = fixed, fixed-exponent = 0, round-mode = places, round-precision = 2] % disable scientific notation
		c
		}
		\textbf{Compound} & \multicolumn{2}{c}{\textbf{Quantity}} & \multicolumn{1}{c}{\textbf{\begin{tabular}[c]{@{}c@{}}Molar mass\\ (\unit{\g\per\mole})\end{tabular}}} & \multicolumn{1}{l}{\textbf{\begin{tabular}[c]{@{}c@{}}Final concentration\\ (for \qty{500}{\ul})\end{tabular}}}         \\*[0.4cm]
		DAPI              & 1                                     & \unit{\ul}                                                                                             & \NA                                                                                                             & 1:500 \\
		\pbs{}            & 499                                   & \unit{\ul}                                                                                             & \NA                                                                                                             & \NA
	\end{tabular}
\end{table}

\clearpage

\section*{\textendash{} Resources \textendash{}}
\markright{Resources}

\begin{enumerate}
	\item \textbf{HCR in Placozoa}
		\begin{itemize}
			\item Najle, S. R., Grau-Bové, X., Elek, A., Navarrete, C., Cianferoni, D., Chiva, C., ... \& Sebé-Pedrós, A. (2023). Stepwise emergence of the neuronal gene expression program in early animal evolution. \textit{Cell}, \textit{186}(21), 4676-4693. \ulhref{https://doi.org/10.1016/j.cell.2023.08.027}{10.1016/j.cell.2023.08.027}
		\end{itemize}
\end{enumerate}

\end{document}


